\chapter{Discussion}\label{ch:discussion}

The preceding chapter provided a systematic, evidence-driven analysis of the operations implemented in NELS. What follows is a more reflective look at that analysis, critically asking what it actually means for the questions this thesis set out to answer. In short, this thesis asked how far latent spaces can be explored systematically without touching the underlying model, which structural and perceptual patterns and instabilities emerge along the way, and where the limits of interpretability and controllability of these latent spaces lie.

A recurring pattern across several operations was that outcomes depended more strongly on the semantic relationship between the inputs than on the geometry of the transformation itself. Interpolation provides a clear example. Although LERP and SLERP follow fundamentally different paths, one linear, the other curved to preserve vector length, they produced nearly indistinguishable images even when their embedding-space trajectories visibly diverged. Trajectory geometry therefore offered little reliable insight into the resulting visual transition. Consequently, source-prompt similarity was a much stronger indicator of whether the change appeared gradual or abrupt. Direction Walk showed the same dependency from another perspective. A direction could be transferred smoothly through embedding space, but it produced a coherent visual edit only when the new starting context was semantically compatible with that direction. Neither the magnitude nor the symmetry of the transfer predicted this outcome. Overall, the observed structural and perceptual behavior was governed less by the mechanics of the operation than by the semantic compatibility between the existing content and the transformation applied to it.

Latent Probe revealed a second, more fundamental limitation: proximity in the compact representation does not necessarily correspond to semantic similarity. Local geometry in the projected space contained enough information to recover a point's embedding and cluster context above chance, supporting the feasibility of systematic exploration without modifying the underlying model. However, reconstructing that position as a prompt often produced entirely different content, even when the reconstructed embedding remained close to the target. Semantic Axis Projection showed a related limitation. A user-defined primary axis generalized reliably to unseen prompts from the same themes, whereas the automatically derived second axis was rarely a clear, interpretable contrast and instead reflected the strongest residual variation after removing informtaion of first axis. Systematic exploration is therefore possible, but primarily as an approximate structural analysis rather than as a precise or self-sufficient reconstruction of meaning.

With Consistency Check, however, this thesis was not able to determine whether the underlying latent space is actually stable during navigation. The operation's advertised continuous variation-strength control behaved, in practice, like a much coarser switch, and its two intended stability signals frequently disagreed with each other rather than confirming the same picture. This does not mean that navigation through this space is itself unstable, only that the specific technique built here to measure it could not reliably show one way or the other.

The third question, concerning the limits of interpretability and controllability of these latent spaces, can already be answered from the patterns described above. Direction Walk and Latent Probe both show that this system's behavior is genuinely difficult to predict and to control from an operation's own parameters alone, whether that means anticipating how a transferred direction will be expressed or assuming that a recovered position also carries the original meaning. What remains unresolved is how strongly ordinary users would perceive these limitations in practice. The qualitative assessments in \hyperref[ch:results]{Chapter~5} reflect the author's interpretation rather than evaluations by independent participants. Thus, a dedicated user study is the logical next step to test whether these automatically measured limits also hold up against real human perception.

These findings should be interpreted strictly within the scope of the evaluation. They are based on a single diffusion model, one text encoder, and a fixed generation configuration, while several experiments used deliberately selected samples rather than randomly sampled cases. The conclusions are therefore supported under the tested conditions but should not be generalized into a broader theory of diffusion latent-space behavior.
